\documentclass[11pt,a4paper]{article}
\usepackage[utf8]{inputenc}
\usepackage[spanish]{babel}
\usepackage[margin=2.0cm]{geometry}
\usepackage{hyperref}
\usepackage{enumitem}

\title{\textbf{Propuesta de Tesis de Licenciatura en Ciencia de Datos}\\
\large Tomografía Sísmica y Análisis de Datos Geológicos}
\author{Director: Dr. Martín D. Maas\\
Departamento de Matemática, FCEyN-UBA\\
Instituto de Matemática Alberto Santaló (IMAS), UBA-CONICET}
\date{}

\begin{document}

\maketitle

\section*{Resumen}

La geología computacional es un área interdisciplinaria donde convergen computación, probabilidades, física, optimización, análisis numérico y machine learning. Esta propuesta se centra en el estudio de métodos para generar imágenes del interior de la Tierra a partir de datos sísmicos (tomografía sísmica), una técnica fundamental tanto en aplicaciones de ciencia básica como en la industria.

\section*{Contexto y Motivación}

La tomografía sísmica es una herramienta clave para comprender la estructura interna de la Tierra. En ciencia básica, se utiliza para desarrollar \emph{modelos globales} que permiten obtener conocimiento geológico fundamental sobre el interior del planeta, que sigue siendo en gran medida desconocido. En aplicaciones industriales, estas técnicas son cruciales para la exploración de recursos naturales.

El área presenta desafíos computacionales y matemáticos interesantes que requieren la integración de múltiples disciplinas: procesamiento de señales, optimización de gran escala, análisis de datos de alta dimensión, y métodos de aprendizaje automático.

\section*{Objetivos}

La tesis puede enfocarse en uno o varios de los siguientes aspectos:

\begin{itemize}[leftmargin=*]
    \item \textbf{Algoritmos de inversión:} Estudio e implementación de algoritmos para resolver el problema inverso de tomografía sísmica, incluyendo métodos de optimización y regularización.
    
    \item \textbf{Análisis de datos no supervisado:} Aplicación de técnicas como Análisis de Componentes Principales (PCA), PCA robusto, y mapas auto-organizados (Self-Organizing Maps) a atributos sísmicos para identificar patrones geológicos.
    
    \item \textbf{Procesamiento de datos reales:} Trabajo con datasets geológicos públicos, particularmente de modelos globales de tomografía sísmica y registros de terremotos.
\end{itemize}

El alcance específico se definirá en función de los intereses del estudiante y puede orientarse más hacia optimización, análisis de datos, o una combinación de ambos.

\section*{Metodología}

El desarrollo de la tesis incluirá:

\begin{enumerate}[leftmargin=*]
    \item Revisión bibliográfica de métodos de tomografía sísmica y técnicas de análisis de datos aplicadas a geofísica.
    
    \item Implementación desde primeros principios de los métodos estudiados, con énfasis en la comprensión de los fundamentos matemáticos y físicos.
    
    \item Trabajo con datos reales o simulados, dependiendo del enfoque específico del proyecto.
    
    \item Análisis de resultados y comparación con métodos existentes en la literatura.
\end{enumerate}

\section*{Qué aprenderá el estudiante}

El estudiante adquirirá conocimientos sobre:

\begin{itemize}[leftmargin=*]
    \item Los tipos de problemas que aparecen en el procesamiento y análisis de datos geológicos, desde la toma de mediciones hasta el análisis final para interpretar estructuras de interés.
    
    \item Formulación matemática de problemas inversos y técnicas de solución computacional.
    
    \item Métodos avanzados de análisis de datos de alta dimensión aplicados a datos científicos reales.
    
    \item Integración de conocimientos de múltiples disciplinas: física de ondas, optimización numérica, estadística y machine learning.
    
    \item Programación científica y buenas prácticas en el desarrollo de software para análisis de datos.
\end{itemize}

\section*{Referencias Relevantes}

\begin{itemize}[leftmargin=*]
    \item Rawlinson, N., Pozgay, S., \& Fishwick, S. (2010). ``Seismic tomography: A window into deep Earth.'' \emph{Physics of the Earth and Planetary Interiors}, 178(3-4), 101--135.
    
    \item Toksoz, M. et al. (2016). ``Geologic pattern recognition from seismic attributes: Principal component analysis and self-organizing maps.'' \emph{Interpretation}, 4(1).
    
    \item Modelos globales de tomografía sísmica: \url{https://www.iris.edu/hq/}
\end{itemize}

\section*{Perfil del Estudiante}

Esta propuesta es ideal para estudiantes de Licenciatura en Ciencia de Datos con interés en:
\begin{itemize}[leftmargin=*]
    \item Aplicaciones científicas del análisis de datos
    \item Optimización y métodos numéricos
    \item Problemas interdisciplinarios
    \item Trabajo con datos reales de gran escala
\end{itemize}

\section*{Proyección}

Esta tesis puede servir como base para estudios de posgrado interdisciplinarios en áreas como computación científica, geofísica computacional, o ciencia de datos aplicada, con potencial para desarrollar investigación con impacto tanto en ciencia básica como en aplicaciones industriales.

\vspace{1cm}

\noindent\textbf{Contacto:} martin.d.maas@gmail.com

\end{document}
